%% DRAFT — extreme-temperature YSZ configuration-survey table, for author
%% review before insertion into the supplementary material (R. Guymon,
%% PR #12, 2026-07-24: "create the YSZ grain growth table you were talking
%% about"). Compiles standalone (`pdflatex ysz_survey_table_draft.tex`).
%%
%% Source of every row: the laboratory records committed at the repo root —
%% "Grain Growth Summary.pdf" (per-run pages with setup, dates, grain sizes,
%% named trace files, and observations) and "Induction Furnace Key
%% Takeaways (1).pdf" (the four-configuration summary slide). Nothing here
%% is inferred beyond those two documents; blank cells mean the record is
%% blank.
%%
%% Intended destination: SI Sec. SIV (high-temperature extension), as
%% \begin{landscape}...\end{landscape} like Table S1. CAUTION on insertion:
%% Sec. SIV precedes the linkage table, so this would become Table S1 and
%% renumber linkage -> S2 and BOM -> S3; the main text's hard-coded
%% "Table S1" mentions must be updated in the same commit.
\documentclass[10pt,letterpaper]{article}

\usepackage[margin=2.0cm]{geometry}
\usepackage{booktabs}
\usepackage{array}
\usepackage{pdflscape}
\usepackage{textcomp}

\newcommand{\dC}{\,$^{\circ}$C}
\newcommand{\um}{\,\textmu m}
\newcommand{\alttext}[1]{\par\smallskip{\footnotesize\emph{Alt text:} #1\par}}
\newcolumntype{L}[1]{>{\raggedright\arraybackslash}p{#1}}

\pagestyle{empty}

\begin{document}

\begin{landscape}
\begin{table}[htbp]
\centering
\caption{Survey of the extreme-temperature YSZ grain-growth configurations
attempted in the induction furnace, in chronological order. Configurations
list the sample assembly top to bottom as loaded in the chamber.
Temperatures are ratio-pyrometer readings and are nominal for the run;
durations are time at temperature. Grain sizes are from optical micrographs
of the recovered specimens. Trace files name the raw heating-profile logs
retained in the laboratory archive; runs without an entry predate
systematic trace retention. The 2500\dC{} / 45\,min row is the successful
anneal reported in the main text (Fig.~7).}
\label{tab:ysz-survey}
\footnotesize
\setlength{\tabcolsep}{4pt}
\renewcommand{\arraystretch}{1.25}
\begin{tabular}{L{1.9cm}L{1.5cm}L{1.3cm}L{4.0cm}L{2.3cm}L{6.6cm}L{2.7cm}}
\toprule
Date & Nominal T & Duration & Configuration (top\,$\rightarrow$\,bottom) &
YSZ grain growth & Outcome & Trace file(s) \\
\midrule
2023-07-04 & 2000\dC & 80\,h &
BN / YSZ / Ta &
--- &
Specimen not usable for grain-size measurement. &
\texttt{2000forWeek.fig}, \texttt{fullTime.xlsx} \\

2024-01-31 & 2000\dC & 24\,h &
Ta (20\,mm) / YSZ / Ta (20\,mm) / alumina rod / Teflon rod &
10\um{} $\rightarrow$ 20\um &
Growth too slow at 2000\dC{}; motivated moving subsequent runs to
2500\dC{}. &
--- \\

2024-03-11 & 2500\dC & 45\,min &
Ta (20\,mm) / YSZ / Ta (20\,mm) / BN crucible / alumina rod / Teflon rod &
20\um{} $\rightarrow$ 90\um &
First successful 2500\dC{} run; clean specimen. Quartz tube softened,
deformed onto the Ta, and was punctured; the interlock stopped the run.
Lesson adopted: fixture the quartz with no-slack pyrometer cables. &
--- \\

2024-04-29 & 2475\dC & 8\,h &
Ta / YSZ / BN &
$\sim$10\um{} $\rightarrow$ $\sim$400\um{} (tentative) &
Stable for the full 8\,h. Ta vapor deposited on the YSZ; the specimen was
polished and thermally etched before grains could be measured, so the
final size is uncertain. Quartz near the Ta softened slightly. &
\texttt{Ta on top.xlsx}, \texttt{2500for8hr.xlsx} \\

2024-05-23 --27 & 2130\dC & 80.4\,h &
Ta (large) / YSZ / Ta (thin) / BN / alumina / Teflon &
--- (specimen unsalvageable) &
Susceptor degraded mid-run (a shell, possibly silica--BN--Ta, formed
around the Ta); temperature fell $\sim$500\dC{} over $\sim$5\,h. Quartz
and BN required replacement. &
\texttt{23May2024.lvm}, \texttt{weekLongHeat.fig} \\

2024-07-22 & 1750\dC & 168\,h &
Graphite crucible with YSZ inside, on BN &
--- (specimen consumed) &
YSZ reacted with the graphite and disappeared (carbothermal attack). The
pyrometer likely under-read this run (the susceptor sat partially outside
its view); the true temperature was probably closer to 2000\dC{}. &
\texttt{22July2024.lvm} \\

2024-08-07/08 & $>$2500\dC & 2\,h &
Graphite crucible &
--- &
Pyrometer misaligned; the rise/drop interlock shut the furnace down.
Graphite demanded excessive generator power above 2500\dC{}, motivating
the switch to a large tantalum susceptor block. &
--- \\

2024-08 & 1800\dC & 115\,h &
Ta susceptor (configuration not recorded) &
--- &
Quartz--Ta reaction products formed (``geode'' growths; no contact
required). &
--- \\
\bottomrule
\end{tabular}
\alttext{Table of eight extreme-temperature YSZ runs listing date, nominal
temperature from 1750 to over 2500 degrees Celsius, duration from 45
minutes to 168 hours, the stacked sample-assembly configuration, measured
grain growth where a specimen survived, the outcome or failure mode, and
the retained raw trace files. Only the 45-minute run at 2500 degrees
Celsius produced a clean specimen with measured growth from 20 to 90
micrometers; longer runs failed by contact reactions between the
susceptor, specimen, and surrounding materials.}
\end{table}
\end{landscape}

\end{document}
